% !TeX spellcheck = en_US
\documentclass[french]{yLectureNote}

\title{Électromagnétisme 2 PS}
\subtitle{Physique}
\author{Paulhenry Saux}
\date{\today}
\yLanguage{Français}
%lionels.calmels@cemes.fr
\professor{M.Brut}
\usepackage{graphicx}%----pour mettre des images
\usepackage[utf8]{inputenc}%---encodage
\usepackage{geometry}%---pour modifier les tailles et mettre a4paper
%\usepackage{awesomebox}%---pour les boites d'exercices, de pbq et de croquis ---d\'esactiv\'e pour les TP de PC
\usepackage{tikz}%---pour deiffner + d\'ependance de chemfig
% \usepackage{tabularx}%---pour dimensionner automatiquement les tableaux avec variable X
\usepackage{awesomebox}%---Pour les boites info, danger et autres
\usepackage{menukeys}%---Pour deiffner les touches de Calculatrice
\usepackage{fancyhdr}%---pour les en-t\^ete personnalis\'ees
\usepackage{blindtext}%---pour les liens
\usepackage{hyperref}%---pour les liens (\`a mettre en dernier)
\usepackage{caption}%---pour la francisation de la l\'egende table vers Tableau
\usepackage{pifont}
\usepackage{array}%---pour les tableaux
\usepackage{lipsum}
\usepackage{yFlatTable}
\usepackage{multicol}
\usepackage{tabularx}
\usepackage{booktabs}
\usepackage{esint}% Pour les oint
\newcommand{\Lim}[1]{\lim\limits_{\substack{#1}}\:}
\renewcommand{\vec}{\overrightarrow}
\newcommand{\N}[0]{\mathbb{N}}
\newcommand{\dd}{\mathrm{d}}
\newcommand{\norm}[1]{||\vec{#1}||}
\newcommand{\fo}{\psi(\vec{r},t)}
\newcommand{\foe}{\psi(\vec{r},t)\*}
\newcommand{\HH}{\hat{H}}
\newcommand{\hb}{\hbar}
\newcommand{\lap}{\nabla^2}
\newcommand{\lapcc}{\frac{\partial^2 }{\partial x^2}+\frac{\partial^2 }{\partial y^2}+\frac{\partial^2 }{\partial z^2}}
\newcommand{\E}{\vec{E}(M)}
\newcommand{\B}{\vec{B}(M)}
\newcommand{\V}{V(M)}
\newcommand{\er}{\vec{e_{\rho}}}
\newcommand{\ep}{\vec{e_{\varphi}}}
\newcommand{\pip}{\pi^+}
\newcommand{\pim}{\pi^-}
\newcommand{\mv}{\mathcal{V}}
\DeclareMathOperator{\Div}{Div}
\newcommand{\rot}{\vec{\mathrm{rot}}}
\newcommand{\grad}{\vec{\mathrm{grad}}}
\setcounter{chapter}{3}
\begin{document}
\chapter{Puissance et énergie du champ Électro-magnétique}
\section{Puissance cédée par le champ à la matière}

\begin{definition}[Puissance volumique de Lorentz]
La puissance fournie par un champ électromagnétique $(\vec{E}, \vec{B})$ à une distribution de charges caractérisée par une densité de courant $\vec{j}$ est donnée par la puissance volumique :
$$P_{vol} = \vec{j} \cdot \vec{E}$$
\end{definition}

\tipsInfo{Interprétation physique}{Seul le champ électrique travaille. La force de Laplace magnétique ($\vec{F}_m = q\vec{v} \wedge \vec{B}$) est toujours orthogonale au déplacement, son travail est donc nul. Le terme $\vec{j} \cdot \vec{E}$ représente le transfert d'énergie du champ vers les porteurs de charge (énergie cinétique, effet Joule).}

\begin{theorem}[Puissance cédée à un volume $V$]
La puissance totale cédée par le champ à la matière contenue dans un volume $V$ est :
$$P_{meca} = \iiint_{V} \vec{j} \cdot \vec{E} \, \mathrm{d}\tau$$
\end{theorem}

\section{Le théorème de Poynting}

\subsection{Démonstration locale}

Partons de l'équation de Maxwell-Ampère : $\vec{j} = \frac{1}{\mu_0} \vec{\nabla} \wedge \vec{B} - \epsilon_0 \frac{\partial \vec{E}}{\partial t}$.
En multipliant par $\vec{E}$ :
$$\vec{j} \cdot \vec{E} = \frac{1}{\mu_0} \vec{E} \cdot (\vec{\nabla} \wedge \vec{B}) - \epsilon_0 \vec{E} \cdot \frac{\partial \vec{E}}{\partial t}$$

Utilisons l'identité vectorielle : $\vec{\nabla} \cdot (\vec{E} \wedge \vec{B}) = \vec{B} \cdot (\vec{\nabla} \wedge \vec{E}) - \vec{E} \cdot (\vec{\nabla} \wedge \vec{B})$.
On en déduit : $\vec{E} \cdot (\vec{\nabla} \wedge \vec{B}) = \vec{B} \cdot (\vec{\nabla} \wedge \vec{E}) - \vec{\nabla} \cdot (\vec{E} \wedge \vec{B})$.

D'après Maxwell-Faraday, $\vec{\nabla} \wedge \vec{E} = -\frac{\partial \vec{B}}{\partial t}$, donc :
$$\vec{j} \cdot \vec{E} = -\frac{1}{\mu_0} \vec{B} \cdot \frac{\partial \vec{B}}{\partial t} - \epsilon_0 \vec{E} \cdot \frac{\partial \vec{E}}{\partial t} - \vec{\nabla} \cdot \left( \frac{\vec{E} \wedge \vec{B}}{\mu_0} \right)$$

En remarquant que $\vec{X} \cdot \frac{\partial \vec{X}}{\partial t} = \frac{1}{2} \frac{\partial \|\vec{X}\|^2}{\partial t}$, on obtient la forme locale.

\begin{theorem}[Théorème de Poynting]
$$-\vec{j} \cdot \vec{E} = \frac{\partial u_{em}}{\partial t} + \vec{\nabla} \cdot \vec{\Pi}$$
Où :
\begin{itemize}
    \item $u_{em} = \frac{\epsilon_0 E^2}{2} + \frac{B^2}{2\mu_0}$ est la \textbf{densité volumique d'énergie EM}.
    \item $\vec{\Pi} = \frac{\vec{E} \wedge \vec{B}}{\mu_0}$ est le \textbf{vecteur de Poynting}.
\end{itemize}
\end{theorem}

\subsection{Signification physique des termes}

Le terme $-\vec{j} \cdot \vec{E}$ est la puissance cédée par la matière au champ.\marginTips{Si elle est positive, la matière agit comme une source (ex: générateur). Si elle est négative, le champ cède de l'énergie à la matière (ex: récepteur, effet Joule).}

\warningInfo{Le vecteur de Poynting}{Il représente la densité de courant d'énergie. Son flux à travers une surface fermée orientée vers l'extérieur représente la puissance sortante du volume par rayonnement ou conduction EM.}

\section{Forme intégrale et conservation}

\begin{proposition}[Bilan intégral]
Pour un volume $V$ fixe délimité par une surface $S$ :
$$\frac{\mathrm{d}U_{em}}{\mathrm{d}t} = - \iint_S \vec{\Pi} \cdot \mathrm{d}\vec{S} - \iiint_V \vec{j} \cdot \vec{E} \, \mathrm{d}\tau$$
\end{proposition}

La variation de l'énergie stockée dans le volume dépend de deux processus :
\begin{enumerate}
    \item Le flux du vecteur de Poynting (échange avec l'extérieur).
    \item Le travail des forces de Lorentz (échange avec la matière).
\end{enumerate}

% \section{Applications aux composants}

% \subsection{Le Condensateur (Stockage Électrique)}
% Pour un condensateur plan de capacité $C = \frac{\epsilon_0 S}{e}$ :
% \begin{itemize}
%     \item Champ électrique : $E = \frac{U}{e}$.
%     \item Énergie stockée : $U_e = \iiint \frac{\epsilon_0 E^2}{2} \mathrm{d}\tau = \frac{\epsilon_0}{2} \left(\frac{U}{e}\right)^2 (S \cdot e) = \frac{1}{2} C U^2$.
% \end{itemize}

% \subsection{Le Solénoïde (Stockage Magnétique)}
% Pour un solénoïde de longueur $l$, de $n$ spires par mètre, d'inductance $L = \mu_0 n^2 S l$ :
% \begin{itemize}
%     \item Champ magnétique : $B = \mu_0 n I$.
%     \item Énergie stockée : $U_m = \iiint \frac{B^2}{2\mu_0} \mathrm{d}\tau = \frac{(\mu_0 n I)^2}{2\mu_0} (S \cdot l) = \frac{1}{2} L I^2$.
% \end{itemize}

% \subsection{La Résistance (Effet Joule)}
% Considérons un conducteur cylindrique de rayon $a$, de longueur $l$ et de conductivité $\gamma$.



% \criticalInfo{Origine de l'énergie Joule}{Dans une résistance, le vecteur de Poynting est dirigé \textbf{vers l'intérieur} du fil : $\vec{\Pi} = -\frac{I^2 \rho}{2\pi^2 a^4 \gamma} \vec{e}_\rho$. Cela signifie que l'énergie thermique dissipée par effet Joule ne vient pas "de l'intérieur" du courant mais est apportée latéralement par le champ électromagnétique environnant.}

% \begin{itemize}
%     \item Puissance entrant par la surface latérale ($\rho = a$) :
%     $$P_{entrante} = \iint \vec{\Pi}_{entrant} \cdot \mathrm{d}\vec{S} = \Pi(a) \cdot (2\pi a l) = \frac{I^2 a}{2\pi^2 a^4 \gamma} \cdot 2\pi a l = \frac{l}{\gamma \pi a^2} I^2 = R I^2$$
% \end{itemize}
% On retrouve exactement la puissance dissipée par effet Joule.

% \tipsInfo{Vérification rapide}{Pour savoir si un système stocke ou dissipe, regardez le signe de $\vec{\nabla} \cdot \vec{\Pi}$. Si le flux est entrant (convergent), le système reçoit de l'énergie du champ.}

\section{Applications aux composants : Approche énergétique}
\subsection{Le Condensateur (Stockage Électrique)}

On considère un condensateur plan (surface $S$, écartement $e$). En négligeant les effets de bord, le champ est uniforme entre les armatures.

\begin{flalign*}
    &\text{Établissement du champ :} && \vec{E} = \frac{\sigma}{\epsilon_0}\vec{e}_z = \frac{Q}{\epsilon_0 S}\vec{e}_z &\\
    &\text{Lien avec la tension :} && U = V(0) - V(e) = \int_{0}^{e} \vec{E} \cdot \mathrm{d}\vec{l} = E \cdot e \implies E = \frac{U}{e} &\\
    &\text{Capacité :} && Q = \frac{\epsilon_0 S}{e} U = CU \quad \text{avec} \quad C = \frac{\epsilon_0 S}{e} &
\end{flalign*}

\begin{proposition}[Énergie du condensateur]
L'énergie est localisée dans le champ électrique entre les plaques.
\begin{flalign*}
    & u_e = \frac{1}{2}\epsilon_0 E^2 = \frac{1}{2}\epsilon_0 \left(\frac{U}{e}\right)^2 &\\
    & U_{em} = \iiint_{V} u_e \mathrm{d}\tau = u_e \cdot (S \cdot e) = \frac{1}{2}\epsilon_0 \frac{U^2}{e^2} S e = \frac{1}{2} \left(\frac{\epsilon_0 S}{e}\right) U^2 = \frac{1}{2}CU^2 &
\end{flalign*}
\end{proposition}

\tipsInfo{Interprétation de Poynting}{Lors de la charge, un champ magnétique induit $\vec{B}_{ind}$ apparaît (courant de déplacement). Le vecteur $\vec{\Pi}$ est alors dirigé de l'extérieur vers l'intérieur de l'espace entre les plaques : le condensateur "engloutit" l'énergie du champ pour la stocker.}

\subsection{Le Solénoïde (Stockage Magnétique)}

Pour un solénoïde infini (ou $l \gg a$), le champ est nul à l'extérieur et uniforme à l'intérieur.

\begin{flalign*}
    &\text{Théorème d'Ampère :} && \oint \vec{B} \cdot \mathrm{d}\vec{l} = \mu_0 I_{enl} \implies B \cdot L = \mu_0 (n \cdot L) I \implies \vec{B} = \mu_0 n I \vec{e}_z &\\
    &\text{Inductance :} && \Phi = N \iint \vec{B} \cdot \mathrm{d}\vec{S} = (nl) \cdot (\mu_0 n I S) = (\mu_0 n^2 S l) I = LI &
\end{flalign*}

\begin{proposition}[Énergie du solénoïde]
L'énergie est stockée sous forme magnétique dans le volume intérieur.
\begin{flalign*}
    & u_m = \frac{B^2}{2\mu_0} = \frac{(\mu_0 n I)^2}{2\mu_0} = \frac{1}{2} \mu_0 n^2 I^2 &\\
    & U_{em} = \iiint_{V} u_m \mathrm{d}\tau = u_m \cdot (S \cdot l) = \frac{1}{2} (\mu_0 n^2 S l) I^2 = \frac{1}{2}LI^2 &
\end{flalign*}
\end{proposition}

\subsection{La Résistance (Effet Joule)}

Considérons un fil cylindrique de rayon $a$, longueur $l$, conductivité $\gamma$, parcouru par un courant $I$.

\begin{flalign*}
    &\text{Champ électrique (Ohm) :} && \vec{j} = \gamma \vec{E} \implies \vec{E} = \frac{I}{\gamma \pi a^2} \vec{e}_z = \frac{V}{l}\vec{e}_z &\\
    &\text{Champ magnétique (Ampère) :} && \text{À la surface } (\rho=a) : \vec{B} = \frac{\mu_0 I}{2\pi a} \vec{e}_\theta &\\
    &\text{Vecteur de Poynting :} && \vec{\Pi} = \frac{\vec{E} \wedge \vec{B}}{\mu_0} = \frac{1}{\mu_0} \left( E \vec{e}_z \wedge B \vec{e}_\theta \right) = - \frac{E B}{\mu_0} \vec{e}_\rho &
\end{flalign*}

En remplaçant les expressions de $E$ et $B$ à la surface ($\rho=a$) :
\begin{flalign*}
    & \vec{\Pi}(a) = - \frac{1}{\mu_0} \left( \frac{V}{l} \right) \left( \frac{\mu_0 I}{2\pi a} \right) \vec{e}_\rho = - \frac{V I}{2\pi a l} \vec{e}_\rho &
\end{flalign*}



\criticalInfo{Origine de l'énergie Joule}{Le calcul du flux sortant à travers la surface latérale $S_L$ donne :
$$\oiint \vec{\Pi} \cdot \mathrm{d}\vec{S} = \vec{\Pi}(a) \cdot (2\pi a l \vec{e}_\rho) = - \frac{V I}{2\pi a l} \cdot 2\pi a l = -VI$$
Le flux est \textbf{négatif}, donc de la puissance \textbf{entre} dans la résistance. L'énergie thermique dissipée par effet Joule ($RI^2$) ne provient pas d'un flux d'énergie circulant \textit{dans} le fil avec les électrons, mais d'une "aspiration" de l'énergie électromagnétique du milieu environnant vers l'intérieur du conducteur.}

\checkInfo{Bilan local}{Dans le cas stationnaire ($\frac{\partial u_{em}}{\partial t} = 0$), le théorème de Poynting devient :
$$\vec{\nabla} \cdot \vec{\Pi} = -\vec{j} \cdot \vec{E}$$
Ici, $\vec{j} \cdot \vec{E} = \gamma E^2 > 0$ (puissance cédée à la matière), donc $\vec{\nabla} \cdot \vec{\Pi} < 0$ : le champ diverge négativement, ce qui confirme que la résistance est un puits d'énergie pour le champ EM.}

\end{document}
